% =====================================================================
%  Appendix B -- the concordance.
%
%  This is the instrument that makes "no content omitted" checkable
%  rather than asserted.  Four tables:
%    B.1  every one of the 128 source headings -> a live \S\ref
%    B.2  every self-check question that moved, and why
%    B.3  every \tag{}-numbered source equation -> a live \eqref
%    B.4  every fenced block in the source -> figure / listing / table
%
%  check_coverage.py asserts against all four.
%  The B.1 row body is generated by that script's companion generator and
%  lives in _concordance_rows.tex; regenerate it, do not hand-edit rows.
% =====================================================================

\section{Concordance with the source document}
\label{app:concordance}

The re-architecture in this report is a change of \emph{order}, not of content:
the source's node-by-node parts were separated into theory
(\S\ref{part:ratetheory}--\S\ref{part:weak}), implementation
(\S\ref{part:evaluator}), verification (\S\ref{part:verification}), policy
(\S\ref{part:guards}) and the open register (\S\ref{part:open}). That makes a
concordance necessary rather than decorative: a source Part's argument can now
land in three places, and this appendix is where you find them.

Everything below is mechanically checked by
\code{study/tier1-report/check\_coverage.py}.

% =====================================================================
\subsection{Headings}
\label{app:headings}

All 128 headings of \code{study/tier1.md}, in source order, with the line at
which each begins and the section(s) of this report that carry it. A heading
with more than one target is one whose material was split across layers.

\begingroup
\small
\begin{longtable}{@{}L{1.7cm} L{6.0cm} r L{4.6cm}@{}}
\toprule
\textbf{Source \S} & \textbf{Heading} & \textbf{Line} & \textbf{Here} \\
\midrule
\endfirsthead
\toprule
\textbf{Source \S} & \textbf{Heading} & \textbf{Line} & \textbf{Here} \\
\midrule
\endhead
\bottomrule
\endfoot
--- & Tier 1 --- Microphysics (title) & 1 & \S\ref{sec:status} \\
--- & Status of this document & 16 & \S\ref{sec:status} \\
--- & How to study each node & 45 & \S\ref{sec:howtostudy} \\
--- & Contents / full section list & 60 & \S\ref{sec:readingmap}, \S\ref{app:sectionmap} \\
--- & Part I --- rate theory & 140 & \S\ref{part:ratetheory} \\
1.1 & From cross-section to <sigma v> & 156 & \S\ref{sec:sigmav} \\
1.1.1 & The reaction rate per unit volume & 158 & \S\ref{sec:rate-per-volume} \\
1.1.2 & Double counting: identical particles & 215 & \S\ref{sec:double-counting} \\
1.1.3 & Why the molar convention & 247 & \S\ref{sec:molar} \\
1.1.4 & The code contract & 302 & \S\ref{sec:code-contract} \\
1.2 & The Coulomb barrier and the Gamow peak & 330 & \S\ref{sec:gamow} \\
1.2.1 & The problem & 332 & \S\ref{sec:barrier} \\
1.2.2 & The Gamow factor and the S-factor & 354 & \S\ref{sec:sfactor} \\
1.2.3 & The saddle point: $E_0$ and $\Delta$ & 385 & \S\ref{sec:saddle} \\
1.2.4 & The numbers in this box & 449 & \S\ref{sec:gamow-numbers} \\
1.2.5 & Neutrons have no Gamow peak & 490 & \S\ref{sec:no-gamow} \\
1.2.6 & The $1/v$ law; the reciprocity theorem & 501 & \S\ref{sec:reciprocity} \\
1.3 & Resonances and Hauser--Feshbach & 554 & \S\ref{sec:hf} \\
1.3.1 & The isolated narrow resonance & 556 & \S\ref{sec:narrowres} \\
1.3.2 & Why Hauser--Feshbach dominates here & 591 & \S\ref{sec:hf-dominates} \\
1.3.3 & Where the uncertainty actually is & 626 & \S\ref{sec:uncertainty} \\
1.3.4 & The $\Gamma/D$ criterion & 648 & \S\ref{sec:gammaD} \\
1.3.5 & Ground-state vs stellar rates (SEF) & 677 & \S\ref{sec:sef} \\
1.4 & Photodisintegration & 738 & \S\ref{sec:photodis} \\
1.4.1 & A rate with a thermal photon bath & 740 & \S\ref{sec:photonbath} \\
1.4.2 & Worked crossover Si28 / S32 & 760 & \S\ref{sec:crossover} \\
1.4.3 & Why ${\sim}3$\,GK & 807 & \S\ref{sec:why3gk} \\
1.4.4 & Two assumptions inside the photon rate & 825 & \S\ref{sec:photon-assumptions} \\
1.5 & What the rates do & 857 & \S\ref{sec:whatratesdo} \\
1.6 & Self-check for Part I & 880 & \S\ref{sec:selfcheck-ratetheory}, \S\ref{sec:selfcheck-evaluator} \\
--- & Part II (S2) --- REACLIB & 899 & \S\ref{part:reaclib} \\
II.1 & The seven-coefficient form, derived & 905 & \S\ref{sec:sevencoeff} \\
II.2 & Sets, chapters, labels, the v flag & 949 & \S\ref{sec:labelprops}, \S\ref{sec:vflag-extraction} \\
II.3 & One rate dissected & 978 & \S\ref{sec:onerate} \\
II.4 & From $\lambda$ to the gross flux & 1028 & \S\ref{sec:lambda-to-R} \\
II.5 & Code: the coefficient tensor & 1060 & \S\ref{sec:coefftensor} \\
II.6 & Oracle: the $5\times10^{-12}$ gate & 1089 & \S\ref{sec:oracle-compile} \\
II.6b & The numerical range of $\lambda$ & 1116 & \S\ref{sec:lambda-range} \\
II.6c & Multiple sets, and fit validity & 1150 & \S\ref{sec:multisets} \\
II.7 & A fit library, not a theory & 1174 & \S\ref{sec:fitnottheory} \\
II.8 & Self-check for S2 & 1189 & \S\ref{sec:selfcheck-reaclib}, \S\ref{sec:selfcheck-evaluator} \\
--- & Part III (S3) --- detailed balance & 1206 & \S\ref{part:db} \\
III.1 & The reverse/forward ratio, derived & 1220 & \S\ref{sec:db-ratio} \\
III.2 & The thermal phase-space factor & 1335 & \S\ref{sec:phasespace} \\
III.3 & Partition functions in the ratio & 1368 & \S\ref{sec:pf-in-ratio} \\
--- & (unnumbered) The $G$ vs $(2J{+}1)G$ trap & 1403 & \S\ref{sec:pf-normalisation} \\
--- & (unnumbered) The pf spline past the table & 1429 & \S\ref{sec:pf-spline} \\
III.4 & Three ways to build a reverse & 1474 & \S\ref{sec:threeways} \\
III.5 & The v-flag dissected & 1494 & \S\ref{sec:vflag} \\
III.6 & $\kappa$ as the numerical DB test & 1561 & \S\ref{sec:kappa-db} \\
III.7 & gh-575, derived from the same algebra & 1595 & \S\ref{sec:gh575} \\
III.8 & Code: \code{db\_reverses.py} and the gate & 1714 & \S\ref{sec:db-replacement}, \S\ref{sec:guard-pf} \\
III.9 & Oracles & 1747 & \S\ref{sec:oracles-db} \\
III.10 & The two $\kappa$ conventions & 1768 & \S\ref{sec:two-kappa} \\
III.11 & Coulomb terms in NSE & 1784 & \S\ref{sec:nse-coulomb} \\
III.12 & $Q(T)$, derived and measured & 1839 & \S\ref{sec:qoft} \\
III.13 & Self-check for S3 & 1935 & \S\ref{sec:selfcheck-db}, \S\ref{sec:selfcheck-plasma}, \S\ref{sec:selfcheck-verification} \\
--- & Part IV (S4) --- screening & 1964 & \S\ref{part:plasma} \\
IV.1 & Why screening exists & 1970 & \S\ref{sec:whyscreening} \\
IV.2 & Weak screening: Debye--Huckel & 1993 & \S\ref{sec:debye} \\
IV.3 & The coupling parameter $\Gamma$ & 2057 & \S\ref{sec:gamma} \\
--- & (unnumbered) The strong-screening limit & 2092 & \S\ref{sec:strongscreening} \\
IV.4 & Where the box sits & 2141 & \S\ref{sec:boxgamma} \\
IV.5 & \code{chugunov\_2007} term by term & 2181 & \S\ref{sec:chugunov}, \S\ref{sec:screening-code} \\
IV.6 & Per-reaction pair construction & 2283 & \S\ref{sec:screening-pairs}, \S\ref{sec:screening-code} \\
IV.7 & The screened-$\kappa$ offset, derived & 2316 & \S\ref{sec:screened-kappa} \\
IV.8 & Code and oracle & 2366 & \S\ref{sec:oracle-screening} \\
IV.9 & Screening of \emph{weak} rates & 2382 & \S\ref{sec:weak-screening} \\
IV.10 & Self-check for S4 & 2442 & \S\ref{sec:selfcheck-plasma}, \S\ref{sec:selfcheck-verification} \\
--- & Part V (S5) --- the weak sector & 2459 & \S\ref{part:weak} \\
V.1 & Why weak reactions are different & 2471 & \S\ref{sec:weak-different} \\
V.2 & From the golden rule to $ft$ values & 2512 & \S\ref{sec:goldenrule} \\
V.3 & Gamow--Teller strength & 2559 & \S\ref{sec:gt} \\
V.4 & EC in a degenerate plasma & 2593 & \S\ref{sec:ec-degenerate} \\
V.4b & Pauli blocking --- the ratchet & 2672 & \S\ref{sec:pauli} \\
V.5 & The $(T,\rho Y_e)$ table & 2721 & \S\ref{sec:weaktable}, \S\ref{sec:tabular-path} \\
V.5b & Interpolation error, derived & 2775 & \S\ref{sec:interp-error} \\
V.6 & The three channels and the ledger & 2843 & \S\ref{sec:ledger}, \S\ref{sec:ledger-enforcement} \\
V.7 & Table families and precedence & 2872 & \S\ref{sec:tablefamilies}, \S\ref{sec:duplicate-links}, \S\ref{sec:guard-reconciled} \\
V.8 & Neutrino losses & 2930 & \S\ref{sec:nulosses} \\
V.9 & Code path summary & 2945 & \S\ref{sec:codepath} \\
V.10 & The two networks differ & 2964 & \S\ref{sec:weak-asymmetry} \\
V.11 & Self-check for S5 & 2992 & \S\ref{sec:selfcheck-weak}, \S\ref{sec:selfcheck-evaluator}, \S\ref{sec:selfcheck-verification} \\
--- & Part VI (S6) --- reconciliation & 3022 & \S\ref{part:verification} \\
VI.1 & Why this node exists at all & 3029 & \S\ref{sec:whyS6} \\
VI.2 & The canonical key & 3049 & \S\ref{sec:canonicalkey} \\
VI.3 & Membership vs values & 3089 & \S\ref{sec:membership} \\
VI.4 & The disposition algebra & 3109 & \S\ref{sec:disposition} \\
VI.5 & The measured value comparison & 3152 & \S\ref{sec:measured-values} \\
VI.6 & Appendix-B as a measurement & 3206 & \S\ref{sec:appendixb} \\
VI.7 & The screening determination & 3232 & \S\ref{sec:screening-determination} \\
VI.8 & The Fortran probe & 3249 & \S\ref{sec:probe} \\
VI.8b & The probe's five modes and the grid & 3267 & \S\ref{sec:probe-modes} \\
VI.8c & Reading the statistics honestly & 3308 & \S\ref{sec:reading-stats} \\
VI.9 & What the reconciliation licenses & 3332 & \S\ref{sec:licenses} \\
VI.10 & Self-check for S6 & 3356 & \S\ref{sec:selfcheck-verification} \\
--- & Part VII --- the guards & 3375 & \S\ref{part:guards}, \S\ref{sec:guard-code}, \S\ref{sec:guard-design} \\
--- & Part VIII --- open prerequisites & 3415 & \S\ref{part:open} \\
VIII.A & Deferred by design & 3422 & \S\ref{sec:deferred} \\
VIII.B0 & Corrections to earlier revisions & 3435 & \S\ref{sec:errata} \\
VIII.B & Closed in this revision & 3460 & \S\ref{sec:closed} \\
VIII.C & Genuine gaps & 3492 & \S\ref{sec:gaps} \\
VIII.C.1 & Rate-uncertainty $\to$ $Y_e$ sensitivity & 3496 & \S\ref{sec:gap-sensitivity} \\
VIII.C.2 & Ground-state vs stellar forwards (SEF) & 3515 & \S\ref{sec:gap-sef} \\
VIII.C.3 & The dropped NU column & 3538 & \S\ref{sec:gap-nu} \\
VIII.C.4 & Imposed $(T,\rho)$ vs feedback & 3556 & \S\ref{sec:gap-feedback} \\
VIII.C.5 & Mass and $Q$-value provenance & 3575 & \S\ref{sec:gap-masses} \\
VIII.C.6 & $e^{\pm}$ pairs at the hot, thin corner & 3593 & \S\ref{sec:gap-pairs} \\
VIII.C.7 & Network-reduction prior art & 3633 & \S\ref{sec:gap-reduction} \\
VIII.C.8 & Isomers & 3649 & \S\ref{sec:gap-isomers} \\
VIII.C.9 & Coulomb-corrected NSE in-box & 3663 & \S\ref{sec:gap-nse} \\
VIII.C.10 & Gate-design statistics & 3681 & \S\ref{sec:gap-gatestats} \\
VIII.D & Common-mode blindness & 3695 & \S\ref{sec:commonmode} \\
VIII.E & Second-pass audit: a missing tier & 3719 & \S\ref{sec:tier4} \\
VIII.E.1 & The model side has no study node & 3725 & \S\ref{sec:gap-modelside} \\
VIII.E.2 & $Q(T)$ --- specified, unimplemented & 3786 & \S\ref{sec:gap-qoft} \\
VIII.E.3 & Solver-independent ground truth & 3795 & \S\ref{sec:gap-thirdsolver} \\
VIII.E.4 & Determinism; float32/float64 & 3810 & \S\ref{sec:gap-determinism} \\
VIII.E.5 & Gradients through the conservation map & 3831 & \S\ref{sec:gap-gradients} \\
VIII.E.6 & The deployment coupling contract & 3853 & \S\ref{sec:gap-coupling} \\
VIII.E.7 & Residual sourcing gaps & 3867 & \S\ref{sec:gap-sourcing} \\
VIII.F & Priority across both passes & 3890 & \S\ref{sec:priority} \\
--- & Where Tier 1 hands off & 3939 & \S\ref{sec:handoff} \\
--- & What Tier 1 upgraded & 3941 & \S\ref{sec:handoff} \\
--- & Threads that carry forward & 3955 & \S\ref{sec:handoff} \\
--- & The three ideas to carry & 3976 & \S\ref{sec:handoff} \\
--- & And one habit, from Part VIII & 3993 & \S\ref{sec:handoff} \\
--- & References & 4009 & \S\ref{sec:references} \\

\end{longtable}
\endgroup

\begin{aside}
Three source headings are unnumbered in the original (the $G$ vs $(2J{+}1)G$
trap and the pf-spline note under \S III.3, and the free-energy derivation under
\S IV.3); they carry ``---'' in the first column. The source's irregular
identifiers --- \S II.6b, \S II.6c, \S V.4b, \S V.5b, \S VI.8b, \S VI.8c and
\S VIII.B0 --- are reproduced as they stand.
\end{aside}

% =====================================================================
\subsection{Relocated self-check questions}
\label{app:selfcheck}

The source carries 6 self-check blocks and 44 questions. All 44 survive; 10
moved, because layering put their material in a different section and a question
should not test material the reader has not yet reached. Nothing was added.

\begin{center}
\begin{tabularx}{\textwidth}{@{}L{2.1cm} Y L{3.4cm}@{}}
\toprule
\textbf{From} & \textbf{Question} & \textbf{To} \\
\midrule
\S 1.6 (4) & prefactor $1/2$, \code{dens\_exp} 1 --- what arity? &
  \S\ref{sec:selfcheck-evaluator} \\
\S II.8 (4) & write \code{lam} symbolically for a chapter-8 rate &
  \S\ref{sec:selfcheck-evaluator} \\
\S II.8 (5) & why \code{set\_starts} rather than a loop &
  \S\ref{sec:selfcheck-evaluator} \\
\S II.8 (6) & what breaks if \code{set\_owner} is unsorted &
  \S\ref{sec:selfcheck-evaluator} \\
\S III.13 (6) & why the $\kappa$-floor screen needs no ground truth &
  \S\ref{sec:selfcheck-verification} \\
\S III.13 (7) & for and against \code{test\_exported\_npz\_trips\_the\_gate} &
  \S\ref{sec:selfcheck-verification} \\
\S III.13 (10) & the spurious floor from \code{use\_coulomb\_corr=True} &
  \S\ref{sec:selfcheck-plasma} \\
\S IV.10 (4) & what bounds the composite-intermediate rule &
  \S\ref{sec:selfcheck-verification} \\
\S V.11 (5) & engine raises vs MESA clips &
  \S\ref{sec:selfcheck-evaluator} \\
\S V.11 (9) & bound the error from $\max|\Delta^2\log_{10}\lambda|$ &
  \S\ref{sec:selfcheck-verification} \\
\S V.11 (10) & the engine reads only column 5 &
  \S\ref{sec:selfcheck-evaluator} \\
\bottomrule
\end{tabularx}
\end{center}

Block-level accounting: the source's 6 blocks become 7, because
\S\ref{part:evaluator} acquires one of its own; \S VI.10's six questions arrive
in \S\ref{sec:selfcheck-verification} unchanged and are joined there by the four
relocated above.

% =====================================================================
\subsection{Numbered equations}
\label{app:equations}

The source numbers 29 equations with \code{\textbackslash tag\{\}}. All 29 are
numbered here and referenced with \code{\textbackslash eqref}; one further
equation (the gh-575 displacement, boxed but untagged in the source) is numbered
so that \S\ref{sec:gh575} and \S\ref{sec:appendixb} can point at the same
object.

\begin{center}
\begin{tabularx}{\textwidth}{@{}L{1.2cm} Y L{1.6cm} @{\hskip 6mm} L{1.2cm} Y L{1.6cm}@{}}
\toprule
\textbf{Src} & \textbf{Content} & \textbf{Here} &
\textbf{Src} & \textbf{Content} & \textbf{Here} \\
\midrule
(1.1)  & \sigv{} from $\sigma(E)$        & \eqref{eq:sigmav} &
(3.4)  & the six-factor DB ratio         & \eqref{eq:dbratio} \\
(1.2)  & gross flux $R_j$                & \eqref{eq:grossflux} &
(3.5)  & the molar constant \fac{}       & \eqref{eq:fac} \\
(1.3)  & the Gamow factor                & \eqref{eq:gamowfactor} &
(3.6)  & mean excitation $\langle E^*\rangle$ & \eqref{eq:meanexc} \\
(1.4)  & the S-factor                    & \eqref{eq:sfactor} &
(3.7)  & $Q_j(T)$                        & \eqref{eq:qoft} \\
(1.5)  & the Gamow peak $E_0$            & \eqref{eq:gamowpeak} &
(4.1)  & Debye--H\"uckel $h_{\rm weak}$  & \eqref{eq:debye} \\
(1.6)  & the width $\Delta$              & \eqref{eq:gamowwidth} &
(4.2)  & $\Gamma_{12}$ and $\ztil$       & \eqref{eq:gammadef} \\
(1.7)  & $\tau$ at the peak              & \eqref{eq:tau} &
(4.3)  & the screened-$\kappa$ offset    & \eqref{eq:screenedkappa} \\
(1.8)  & the non-resonant rate           & \eqref{eq:nonresonant} &
(4.4)  & the plasma state                & \eqref{eq:plasmastate} \\
(1.9)  & $\tau$ in practical units       & \eqref{eq:taupractical} &
(4.5)  & the ion plasma temperature      & \eqref{eq:plasmatemp} \\
(1.10) & the narrow-resonance rate       & \eqref{eq:narrowres} &
(4.6)  & the quantum-corrected \gtil{}   & \eqref{eq:gammatilde} \\
(1.11) & reciprocity                     & \eqref{eq:reciprocity} &
(4.7)  & \code{chugunov\_2007} $h$       & \eqref{eq:chugunov} \\
(2.1)  & the REACLIB seven-coefficient form & \eqref{eq:reaclib} &
(5.1)  & the allowed $\beta$ rate        & \eqref{eq:betarate} \\
(3.1)  & Saha number density             & \eqref{eq:saha} &
(5.2)  & the phase-space integral $f$    & \eqref{eq:phasespaceint} \\
(3.2)  & pairwise balance                & \eqref{eq:dbbalance} &
(5.3)  & the EC integral                 & \eqref{eq:ecintegral} \\
(3.3)  & the density ratio               & \eqref{eq:dbdensities} &
(new)  & the gh-575 displacement         & \eqref{eq:gh575} \\
\bottomrule
\end{tabularx}
\end{center}

% =====================================================================
\subsection{Diagrams, listings and dumps}
\label{app:diagrams}

The source carries 34 fenced blocks. None is dropped: four became TikZ figures,
the rest are set verbatim in the section that now owns them. (Two further
blocks, indented inside list items at source lines 293--297 and 1730--1733, are
also carried --- into \S\ref{sec:code-contract} and \S\ref{sec:db-replacement}
respectively.)

\begingroup
\small
\begin{longtable}{@{}r L{6.4cm} L{2.4cm} L{3.4cm}@{}}
\toprule
\textbf{Lines} & \textbf{Block} & \textbf{Rendered as} & \textbf{Here} \\
\midrule
\endfirsthead
\toprule
\textbf{Lines} & \textbf{Block} & \textbf{Rendered as} & \textbf{Here} \\
\midrule
\endhead
\bottomrule
\endfoot
47--52 & the four-step study loop & TikZ figure & Fig.~\ref{fig:studyloop} \\
76--134 & the full section list & listing & \S\ref{app:sectionmap} \\
315--321 & \code{engine.evaluate\_fluxes} assembly & listing &
  \S\ref{sec:code-contract} \\
862--869 & the $\alpha$ ladder, annotated & TikZ figure &
  Fig.~\ref{fig:alphaladder} \\
965--968 & \code{vflag\_sets} extraction $+$ the gate call & listing &
  \S\ref{sec:vflag-extraction} \\
983--991 & \code{si28(a,g)s32} and its v-flag reverse & listing &
  \S\ref{sec:onerate} \\
1035--1041 & \code{evaluate\_lambda}: set sum then corrections & listing &
  \S\ref{sec:lambda-to-R} \\
1154--1157 & \code{c12(a,g)o16}'s two sets & listing &
  \S\ref{sec:multisets} \\
1417--1421 & spins and $\ln G$ carried separately & listing &
  \S\ref{sec:pf-normalisation} \\
1433--1436 & the pf spline, \code{ext="const"} & listing &
  \S\ref{sec:pf-spline} \\
1460--1465 & the pf incidence signs & listing & \S\ref{sec:pf-spline} \\
1554--1557 & \code{PfGateError} & listing & \S\ref{sec:guard-pf} \\
1601--1607 & MESA's \code{No == 1} branch (Fortran) & listing &
  \S\ref{sec:gh575} \\
1716--1718 & \code{replace\_vflag\_reverses} signature & listing &
  \S\ref{sec:db-replacement} \\
2083--2087 & \code{gamma\_e\_fac} and $\Gamma$ & listing & \S\ref{sec:gamma} \\
2290--2293 & \code{screening\_pairs} $\to$ \code{screen\_map} & listing &
  \S\ref{sec:screening-code} \\
2310--2312 & \code{SCREENING\_ALLOWED} & listing &
  \S\ref{sec:guard-screening} \\
2394--2396 & the weak table's eight columns & listing &
  \S\ref{sec:weak-screening} \\
2411--2413 & the engine reads only column 5 & listing &
  \S\ref{sec:weak-screening} \\
2501--2506 & three places weak columns stay unpaired & listing &
  \S\ref{sec:weak-different} \\
2725--2731 & \code{Ni56\_to\_Co56\_weaktab} shape and grids & listing &
  \S\ref{sec:weaktable} \\
2756--2760 & clamped searchsorted $+$ bilinear & listing &
  \S\ref{sec:tabular-path} \\
2766--2769 & \Ye{} from the composition & listing &
  \S\ref{sec:tabular-path} \\
2857--2861 & the charge-to-lepton closure assertion & listing &
  \S\ref{sec:ledger-enforcement} \\
2892--2894 & \code{DEFAULT\_TABULAR\_ORDERING} & listing &
  \S\ref{sec:guard-reconciled} \\
2906--2909 & the ordering guard & listing & \S\ref{sec:guard-reconciled} \\
2920--2924 & duplicate-link resolution & listing &
  \S\ref{sec:duplicate-links} \\
2947--2958 & the weak-sector code path & TikZ figure &
  Fig.~\ref{fig:codepath} \\
3054--3057 & the canonical key grammar & listing &
  \S\ref{sec:canonicalkey} \\
3098--3102 & \code{disposition=None} on purpose & listing &
  \S\ref{sec:membership} \\
3122--3126 & the MESA\_ONLY hard failure & listing &
  \S\ref{sec:disposition} \\
3269--3276 & the probe's five modes & listing & \S\ref{sec:probe-modes} \\
3293--3298 & \code{crosscheck/grids.py} & listing & \S\ref{sec:probe-modes} \\
3746--3772 & the missing Tier 4 & TikZ figure & Fig.~\ref{fig:tier4} \\
\end{longtable}
\endgroup

% =====================================================================
\subsection{What is deliberately \emph{not} in this report}
\label{app:notincluded}

Nothing from the source. Two classes of material are referenced but not
reproduced, and both are referenced in the source the same way:

\begin{itemize}
\item \textbf{Tier~0.} Cross-references of the form \S0.x, \S I.x, \S V.x point
      into \code{study/tier0.md} and its report. They are pointers in the
      source and remain pointers here.
\item \textbf{The audit log.} \code{study/tier1-audit.md} records, claim by
      claim, which passage of which source supports which statement, including
      every correction made during the 2026-08-13 pass with its before\,/\,after.
      \S\ref{sec:errata} carries the corrections themselves; the supporting
      passages stay in the log.
\end{itemize}
